% Generated draft: requires mathematical and source review.
\documentclass[11pt]{article}
\usepackage[margin=25mm]{geometry}
\usepackage{fontspec}
\setmainfont{Latin Modern Roman}
\usepackage{amsmath,amssymb,mathtools,unicode-math}
\setmathfont{Latin Modern Math}
\usepackage{enumitem,needspace,xurl,hyperref}
\setlength{\parindent}{0pt}
\setlength{\parskip}{5pt}
\setlength{\emergencystretch}{3em}
\newcommand{\N}{\mathbb{N}}
\newcommand{\Z}{\mathbb{Z}}
\newcommand{\Q}{\mathbb{Q}}
\newcommand{\R}{\mathbb{R}}
\newcommand{\C}{\mathbb{C}}
\newcommand{\F}{\mathbb{F}}
\newcommand{\A}{\mathbb{A}}
\newcommand{\PP}{\mathbb{P}}
\newcommand{\E}{\mathbb{E}}
\newcommand{\eps}{\varepsilon}
\newcommand{\dd}{\,\mathrm{d}}
\begin{document}
\section*{UnsolvedMath ID: 20000227}
\textbf{Source problem code:} catalogue entry
\subsection{Definitions and mathematical statement}
The complete original text supplied for this ID is only the heading: ``Problem 7. (Suggested by Omid Amini) Here is a `weak version' of Nagata's conjecture.'' No actual conjectural inequality or non-bigness assertion follows in the record. Although the verification field says inspection of page 4 of the six-page workshop PDF recovered a complete block, that block is not included. An exact original statement therefore cannot be recovered from the supplied evidence alone.

The attached report adopts the following blowup interpretation. Let \(S\to\PP^2\) be the blowup of \(s\) distinct points, \(L\) the pullback of a line, and \(E_1,\ldots,E_s\) the exceptional curves. For an integral class \(D=dL-\sum_i m_iE_i\), define its Nagata excess by
\[
\epsilon(D)=\sum_{i=1}^s m_i-\sqrt{s}\,d.
\]
The report uses ``standard integral class'' without defining that term; nonnegative degree/multiplicity conditions, any ordering or Cremona-reduction condition, the field, and the range of \(s\) must be fixed before this becomes a precise original conjecture. Its reduction asserts, for each fixed point configuration and the intended class domain,
\[
\exists D\text{ effective with }\epsilon(D)>0
\quad\Longleftrightarrow\quad
\exists B\text{ big with }\epsilon(B)\geq0.
\]
For a surface, bigness means positive volume
\(\operatorname{vol}(B)=\limsup_{n\to\infty}2h^0(S,\mathcal O_S(nB))/n^2>0\).
Thus the inferred weak assertion is absence of big classes in the closed forbidden halfspace, while the equivalent effective assertion excludes strict positive excess. This is the report's reconstruction, not wording present in the original heading.

The reported forward conversion is
\[
B_k=kD+(s+1)L-\sum_{i=1}^sE_i,
\]
using the ample added class \(A=(s+1)L-\sum_iE_i\). It claims a quantitative integer threshold making \(\epsilon(B_k)\geq0\), and the lower bound
\[
\operatorname{vol}(B_k)\geq s^2+s+1.
\]
The summary says the threshold is exact but does not print it. In the indicated notation the excess condition itself is
\(k\epsilon(D)+s-(s+1)\sqrt{s}\geq0\); no omitted threshold theorem is supplied as a quotation. In the reverse direction, the report asserts that \(nB-L\) is effective for every sufficiently large integer \(n\), by a restriction sequence, and has strict positive Nagata excess when \(\epsilon(B)\geq0\).

The report concludes that this weak non-bigness problem is equivalent to the full classical Nagata prediction for nonsquare \(s\), while retaining a genuine equality-boundary caveat for square \(s\). In the usual Nagata setting the point configuration is very general and \(s\) is in the conjectural range; neither qualification is supplied in the original heading. For nonsquare \(s\), a nonzero integral class of the usual positive-degree type cannot lie on \(\sum m_i=d\sqrt{s}\); square \(s\) allows such boundary classes. The report does not identify absence of big boundary classes with absence of all effective boundary classes.
\subsection{Short English statement}
Only the heading of Amini's weak Nagata problem is supplied. The report reconstructs a non-bigness question and proposes an equivalence with strict effective Nagata violations, with an explicit class-changing conversion and a square-case boundary qualification.
\subsection{Sources}
\begin{itemize}
\item AIM workshop ``Degenerations in algebraic geometry,'' Problem 7, suggested by Omid Amini: \url{https://aimath.org/pastworkshops/degenalggeomproblems.pdf}. The supplied verification refers to page 4 of a six-page PDF but does not include the purported recovered block; the PDF was not accessed here.
\item Local canonical registry \texttt{aim-algebraic-geometry-notes.json}, \texttt{notes[226]}, tag \texttt{problem}; permanent UnsolvedMath ID \(20000227\), source code \texttt{catalogue entry}. The actual original and clean statements contain only the heading and Amini attribution. Original remarks and literature field are empty.
\item Supplied research attempt 1 proposing the effective/big equivalence and amplification construction. No full authoritative citation, complete original conjecture, or definition of standard class is supplied. The report is machine generated, not peer reviewed, and requires independent verification before citation.
\end{itemize}
\end{document}
